\documentclass[12pt]{article}
\usepackage[utf8]{inputenc}
\usepackage[T1]{fontenc}
\usepackage{hyperref}
\title{Tai Daeng Language Resources for K-12 Curriculum Development}
\author{Compiled by Claude Sonnet 5 (claude-sonnet-5) for Andr\'as Kornai}
\date{}
\begin{document}
\maketitle

\section{Introduction}
Tai Daeng (Red Tai; ISO 639-3 tyr, Glottocode taid1249) is a Tai language spoken in northwestern Vietnam and northeastern Laos. Search results give about 80,000 native speakers within an ethnic population of roughly 100,000. Many speakers in Vietnam identify as Tai Dam (Black Tai) instead. The searches found reference and background pages but no school materials.

\section{Primary Educational Resources}
No primer, textbook or school programme for Tai Daeng turned up.

\section{Linguistic and Reference Materials}
\begin{itemize}
\item Tai Daeng language (Wikipedia): \url{https://en.wikipedia.org/wiki/Tai_Daeng_language}
\item Tai Viet script (Wikipedia). This is the script associated with Tai Daeng: \url{https://en.wikipedia.org/wiki/Tai_Viet_script}
\item Omniglot page on the related Tai Dam script: \url{https://mail.omniglot.com/writing/taidon.htm}
\item Cocoon (Huma-Num) catalogue entry for tyr. It lists documentation types including New Testament, texts, phonology and a dictionary item attributed to Hudak (1994). I could not confirm the title of that work: \url{https://cocoon.huma-num.fr/exist/crdo/lexvo.xql?code=tyr}
\item Tai Daeng numerals page, Max Planck Institute for Evolutionary Anthropology: \url{https://lingweb.eva.mpg.de/channumerals/Tai-Daeng.htm}
\end{itemize}

\section{Religious and Community Resources}
\begin{itemize}
\item Joshua Project profile for Tai Daeng in Laos. It reports that Bible translation has started: \url{https://joshuaproject.net/index.php/people_groups/15186/LA}
\end{itemize}

\section{Digital and Audio Resources}
\begin{itemize}
\item Verbix endangered-languages page on Tai Daeng: \url{https://docs.verbix.com/EndangeredLanguages/TaiDaeng}
\end{itemize}
I found no audio or video material.

\section{Curriculum Development Context}
The language is spoken across two countries and has a related script, Tai Viet, which has a Unicode block. The sources found here do not describe school use. Since many Vietnamese speakers identify as Tai Dam, curriculum planning would need to work with communities in both countries and to settle which variety and script to use. I could not confirm the Hudak dictionary.
\end{document}